%%%PAGE 105 rot=0 legibility=good summary=光的偏振（自然光/偏振光对比、应用）；麦克斯韦电磁场理论；电磁波的发射与接收（调幅调频、调谐、检波）；狭义相对论（相对性原理、光速不变、质速关系、高速效应）
%%%BLOCK id=105-01 ch=15 sec=光的偏振 kind=knowledge alt=none cont=none y=0.02-0.24
①\quad \textbf{光的偏振}

\begin{tabular}{@{}l p{0.36\linewidth} p{0.36\linewidth}@{}}
 & 自然光 & 偏振光 \\
来源 & 从普通光源发出的光 & 自然光通过起偏器后的光 \\
 & 在垂直于光的传播方向的平面内，光振动沿任意方向，且沿各方向振动的光强度相同 & 在垂直于光的传播方向的平面内. 光振动沿特定方向，\newline 应用：加偏振滤光片的照相机\unclear{镜头}\newline 液晶显示器、立体电影、\newline 消除车灯眩光 \\
\end{tabular}

%FIG p105_a 0.17 0.118 0.455 0.178
%FIG p105_b 0.47 0.1325 0.545 0.158
\noindent\begin{minipage}[c]{0.6\linewidth}\notefig[1.0]{figures/p105_a.png}\end{minipage}\begin{minipage}[c]{0.2\linewidth}\notefig[1.0]{figures/p105_b.png}\end{minipage}

横波\quad 通过偏振片、反射、折射变成只沿某方向振动的光

$V_{\text{传播}}\perp V_{\text{振动}}$
%%%END
%%%BLOCK id=105-02 ch=14 sec=麦克斯韦电磁场理论 kind=knowledge alt=none cont=none y=0.24-0.34
②\quad \textbf{麦克斯韦电磁场理论}

③\quad 变化电/磁场在周围产生变化的磁/电场 % 圈号②③位于左页边，对应关系按垂直位置推定
%%%END
%%%BLOCK id=105-03 ch=14 sec=电磁波的发射与接收 kind=knowledge alt=none cont=none y=0.35-0.53
④\quad \textbf{电磁波的发射与接收，}\quad 电视、雷达、移动电话
\begin{itemize}
\item 发射电磁波要开放的高频振荡电路.\\ 并对电磁波根据信号的强弱进行\underline{调幅}、\underline{调频} % 原稿此两行用左大括号括在一起
\item 接收电磁波需要能产生电谐振的调谐电路，再把信号从高频电流中解调出来\\ 调幅波的解调也叫\underline{检波} % 原稿此两行用左大括号括在一起
\end{itemize}
%%%END
%%%BLOCK id=105-04 ch=15 sec=狭义相对论 kind=knowledge alt=none cont=none y=0.54-0.88
⑤\quad \textbf{相对论.}
\begin{itemize}
\item 狭义相对论
 \begin{itemize}
 \item 狭义相对性原理：在不同的惯性系中一切物理规律相同. % 左页边此行旁有一"L"形短线，疑为层级连线
 \item 光速不变原理\quad 真空中的光速在不同惯性参考系中都相同，与光源、观测者间相对运动无关.
 \end{itemize}
\end{itemize}
\[ \left\{\begin{array}{l} m=\dfrac{m_0}{\sqrt{1-\left(\dfrac{v}{c}\right)^{2}}} \\[3ex] E=mc^{2} \end{array}\right. \]

高速下\quad 尺缩\quad 钟慢\quad 质增 % 其后有涂改带遮盖，字迹不可辨
%%%END
%%%PAGEEND 105
